\section{Conclusion}
\label{sec:conclusion}

We built a benchmark in which an agent has to buy its own data, and then ran 59
agents through it on a single model.

Giving an agent a proven library made it worse. From-scratch agents beat
library-assisted ones in all four substrates, by a pooled median \relltwo{} of
$0.137$ against $0.229$, and the deterministic analysis of what they wrote says
why: given a library, all four settle on its default model and its named
acquisition strategies; given nothing, all four converge on a different family
and two of them beat the library's own score. Holding the model fixed, which
framework wrapped it mattered more than how much help it was given. And the
pipeline that uses no language model at all, at a matched budget and no cost,
still beat five of the eight agent cells.

The finding we would most like to see travel is the one about verification.
Thirty-one per cent of scored runs passed every machine-checkable gate while
being physically invalid, and the single clearest example scored better than
most of the corpus on the error metric while inventing a magnetic field in a
vacuum. The runs that misreported their own performance were largely not the
same runs, so neither measurement detects the other.

A benchmark that checks only whether the required files appeared will certify
work that is physically wrong and self-reported dishonestly, and it will do so
silently. Separating \emph{did it deliver}, \emph{is it valid} and \emph{did it
tell the truth} costs three diagnostics computed from predictions the benchmark
already holds. Asking the fourth question --- \emph{did it build what it said it
built} --- costs a parser. Between them they are the difference between
measuring what an agent can do and measuring whether it filed the right
paperwork.
